\section{推荐阅读}

\begin{readingbox}
\textbf{基础}
\begin{itemize}[nosep]
  \item Eugene Izhikevich, \emph{Dynamical Systems in Neuroscience}，用于获得有纪律的状态空间与稳定性语言。
  \item Michael Brin and Garrett Stuck, \emph{Introduction to Dynamical Systems}，用于理解状态演化与局部结构的数学背景。
  \item Steven Strogatz, \emph{Nonlinear Dynamics and Chaos}，用于理解阈值、稳定性与机制变更的直觉，同时避免夸大实证拟合。
\end{itemize}
\textbf{现代研究与综述}
\begin{itemize}[nosep]
  \item 关于大脑中 attractor 与 integrator 网络的综述文献 (Nature Reviews Neuroscience, 2022, medium)，说明为什么状态持久性与方向信息保留是严肃的建模思路。
  \item \emph{Metastability demystified: the foundational past, the pragmatic present and the promising future} (Nature Reviews Neuroscience, 2024, medium)，说明为什么“稳定但可修正”的状态描述往往比粗糙的刚性/流动二分更合适。
  \item \emph{Cortical connectivity, local dynamics and stability correlates of global conscious states} (Communications Biology, 2025, strong，作为生物学正当性来源，而不是本章评分公式的直接验证)。
\end{itemize}
\textbf{阅读纪律}
\begin{itemize}[nosep]
  \item 用教材来为建模语言提供正当性。
  \item 用综述与生物学来源来说明：严肃对待状态、稳定性与转变结构是合理的。
  \item 不要把任何一组来源读成“本章启发式评分已经完成实证校准”的直接证明。
  \item 如果后续章节引入更强的公式，它们必须明确说明这些公式仍然是启发式的、只是单调性的，还是已经获得了超出本章基础的支持。
\end{itemize}
\end{readingbox}